Multiplication by a bounded potential requires more than weak convergence of the unweighted products. A circular auxiliary variable recovers the weight from an ordinary mixed product, including at top degree.

\begin{lemma}\label{lem:weighted-decreasing-products}\index{Bedford--Taylor product!weighted decreasing limits}
Let $a_j\downarrow a$ and $b_{i,j}\downarrow b_i$, $1\leq i\leq r\leq n$, be locally uniformly bounded plurisubharmonic functions on a domain in $\C^n$. Then
\[
 a_j\,\ddc b_{1,j}\wedge\cdots\wedge\ddc b_{r,j}
 \longrightarrow a\,\ddc b_1\wedge\cdots\wedge\ddc b_r
\]
weakly as currents. The assertion also holds for quasi-plurisubharmonic functions after fixed smooth corrections, and after wedging with smooth forms converging locally in $C^\infty$.
\end{lemma}
\begin{proof}
First let $r=n$ and write $\mu_j=\bigwedge_i\ddc b_{i,j}$ and $\mu=\bigwedge_i\ddc b_i$. Work over a relatively compact ball on which $L\leq a_j,a\leq U$. On the product with a disk of radius $R$, where $\log R^2>U+1$, set
\[
 v_j(z,w)=\max\{a_j(z),\log|w|^2\},\qquad
 v(z,w)=\max\{a(z),\log|w|^2\}.
\]
These are locally bounded plurisubharmonic functions, also at $w=0$, and $v_j\downarrow v$. Pull the base potentials to the product. By \cref{thm:bt-input} in dimension $n+1$,
\[
 T_j=\mu_j\wedge\ddc v_j\longrightarrow T=\mu\wedge\ddc v.
\]
The bounded recursion gives $T=\ddc(v\mu)$. Since $\mu$ has top degree in the base, only derivatives in $w$ survive. Distributional testing and Fubini therefore identify $T$ with $\mu(\mathrm dz)$ times the circle measure
\[
 \ddc_w\max\{a(z),\log|w|^2\}
   =\sigma_{\exp(a(z)/2)}.
\]
Here $\sigma_t$ is uniform probability on $|w|=t$: the radial derivative jumps from zero to $2/t$, its Euclidean Laplacian has total mass $4\pi$, and our $\ddc$ divides that mass by $4\pi$. This distributional computation works for bounded Borel $a$; it assumes no continuity in $z$. The identical description holds for $T_j$.

Choose a smooth radial compactly supported function $\vartheta(w)$, supported away from zero and from $|w|=R$, which equals $\log|w|^2$ on all circles with $L\leq\log|w|^2\leq U$. For every compactly supported smooth base test $\xi$,
\[
 \langle T_j,\xi\vartheta\rangle=\int\xi a_j\,\mathrm d\mu_j,
 \qquad
 \langle T,\xi\vartheta\rangle=\int\xi a\,\mathrm d\mu.
\]
Their convergence proves the top-degree assertion.

For $r<n$, wedge with $n-r$ constant positive forms, each a Hessian of a positive quadratic function, and apply the top-degree case. Polarization recovers arbitrary constant test forms: the decomposable forms obtained from positive Hermitian forms span the real forms of bidegree $(n-r,n-r)$. Smooth scalar tests then recover every current coefficient. To treat quasi-plurisubharmonic functions, add a fixed smooth correction to each decreasing sequence, apply the result to the corrected functions, and expand. Subtracting the smooth correction to $a_j$ uses the ordinary decreasing convergence. Finally the local Chern--Levine--Nirenberg bounds give uniformly bounded variation for these finitely many products on compact sets. They allow locally converging smooth coefficient forms to replace fixed ones in the tests.
\end{proof}
